\subsubsection{Markov Kernels from Structural Causal Models}

In this section we will see how Markov kernels appear naturally in the theory of SCMs.

\begin{Def}\label{def:scm_kernels}
  Let $\C{M} = \langle \Sin, \Send, \Srnd, \CX, P, f \rangle$ be an SCM with solution
  $(X_{V\cup W}^{\doit(\xJ)})_{\xJ\in\CXJ}$ (see Definition~\ref{def:scm_solutions}). Then
%  provides an \emph{observational Markov kernel of $\C{M}$} if the mapping
  $$K : \Bcal_{\CXV \times \CXW} \times \CXJ \to [0,1], \quad (D,\xJ) \mapsto P((\XV^{\doit(\xJ)},\XW^{\doit(\xJ)})\in D)$$
  is a Markov kernel $\CXJ \dshto \CXV \times \CXW$. We call it \emph{a Markov kernel of $\C{M}$}, and
  also write it as $P(X_V,X_W \mid \doit(X_J))$. Note that it may not be unique.
  %This is the case if for each $D \in \Bcal_{\CXV \times \CXW}$ the mapping:
  %      \[ \CXJ \to [0,1],\quad \xJ \mapsto P(D|\xJ) \]
  %is measurable. This is guaranteed to be the case if the random field $(\XV^{\xJ},\XW)_{\xJ\in\CXJ}$ is measurable.

%  We will denote such an observational Markov kernel as
%  $P(\XV^{\XJ},\XW \mid \XJ)$ or simply as $P(\XV,\XW \mid \XJ)$.
  If only a subset $O \subseteq V \cup W$ is observed, we also call the marginal Markov kernel
  $\CXJ \dshto \C{X}_O$, also written as $P_{\C{M}}(X_O \mid \doit(X_J))$, an \emph{(observed) Markov kernel of $\C{M}$}.
%  When $\XW$ is latent, we will also call its marginal
%  $P(\XV^{\XJ} \mid \XJ)$ an observational Markov kernel of $\C{M}$.
%  For a hard intervention with target $T \subseteq \Sin \cup \Send \cup \Srnd$, we will call an
%  observational Markov kernel of the intervened SCM (of one of the three forms
%  $\C{M}_{\doit(T)}$ or $\C{M}_{\doit(\XT=\xT)}$ or $\C{M}_{\doit(\XT\sim Q_T)}$) 
%  an \emph{interventional Markov kernel of $\C{M}$ for that intervention}.
\end{Def}
  Note that if an SCM has multiple solutions, it has multiple Markov kernels, and the notation
  ``$P(X_V,X_W \mid \doit(X_J))$'' is ambiguous since it does not specify which solution the kernel comes
  from. Later, we will mostly restrict attention to simple SCMs for which the Markov kernel turns out to be unique.
  We can construct Markov kernels by pushing the exogenous distribution through a solution function.
\begin{Rem}
  Given an SCM $\C{M} = \langle \Sin, \Send, \Srnd, \CX, P, f \rangle$
  and a solution function $g : \CXJ \times \CXW \to \CXV$, a Markov kernel of $\C{M}$ is obtained as
  the push-forward
  $$P(\XV,\XW \mid \doit(\XJ)) = (g,\id_{\CXW})_* P(\XW \mid \XJ)$$
  of the Markov kernel
  $$P(\XW \mid \doit(\XJ)) := \bigotimes_{w\in\Srnd} P(X_w),$$
  where we interpret the exogenous distribution specified by $\C{M}$ as a constant Markov kernel $\CXJ \dshto \CXW$.
%
%  Similarly, interventional Markov kernels can be obtained by pushing forward the exogenous distribution through a solution function of the intervened SCM.
\end{Rem}
Again, not all Markov kernels of an SCM can be obtained in this way in case the SCM admits multiple solution functions (see also Example~\ref{eg:not_each_solution_from_push_forward}).

\Joris{TODO: explain this in terms of a sampling algorithm.}

By first intervening on the SCM, constructing a Markov kernel for the intervened SCM, and then conditioning and marginalizing it, the SCM leads to a rich family of Markov kernels that forms a nice playground for causal reasoning, inference, discovery and learning.
